\chapter{Vecteurs, matrices et normes}
\label{chap:vecteurs}

\begin{quote}\itshape
Avant de faire apprendre quoi que ce soit à une machine, apprivoisons les
objets avec lesquels elle travaille : les vecteurs et les matrices. Ce
chapitre pose le décor — en douceur.
\end{quote}

\begin{prerequis}
Savoir additionner et multiplier des nombres, et savoir ce qu'est une
racine carrée. Rien d'autre. Si le symbole $\sum$ vous inquiète, ne vous
arrêtez pas : il est expliqué la première fois qu'il apparaît.
\end{prerequis}

\section{Un vecteur, c'est juste une liste de nombres}

\begin{enclair}
Un \emph{vecteur} de taille $n$, c'est une liste ordonnée de $n$ nombres.
Rien de plus. Par exemple, la position d'un point sur une carte
$(x, y)$ est un vecteur de taille~2. Les pixels d'une image, mis bout à
bout, forment un vecteur (très long).
\end{enclair}

\begin{enia}
En intelligence artificielle, \textbf{tout} devient un vecteur : une
photo, une phrase, un client, une molécule. C'est la première chose que
fait n'importe quel système d'IA — transformer le monde en listes de
nombres, parce que c'est la seule chose qu'une machine sache calculer.
Le vecteur qui représente un objet s'appelle son \emph{plongement}
(\emph{embedding} en anglais).
\end{enia}

\begin{graphe}{Un jeu de données, c'est un nuage de vecteurs. Ici, chaque
fleur est décrite par deux nombres seulement — donc par un vecteur de
$\R^2$ —, et l'espèce se lit presque à l'œil nu.}[nuage-iris]
\begin{axis}[mbaxis, height=5.6cm,
             xlabel={Longueur du pétale (cm)},
             ylabel={Largeur du pétale (cm)},
             xmin=0, xmax=7.5, ymin=0, ymax=3,
             legend pos=north west]
  \addplot[only marks, mark=*,       mark size=1.5pt, mbBlue] coordinates {
    (1.4,0.2) (1.4,0.2) (1.3,0.2) (1.5,0.2) (1.4,0.2) (1.7,0.4)
    (1.4,0.3) (1.5,0.2) (1.4,0.2) (1.5,0.1) (1.5,0.2) (1.6,0.2)};
  \addplot[only marks, mark=triangle*, mark size=2pt,  mbAmber] coordinates {
    (4.7,1.4) (4.5,1.5) (4.9,1.5) (4.0,1.3) (4.6,1.5) (4.5,1.3)
    (4.7,1.6) (3.3,1.0) (4.6,1.3) (3.9,1.4) (4.2,1.5) (4.0,1.0)};
  \addplot[only marks, mark=square*, mark size=1.5pt, mbTeal] coordinates {
    (6.0,2.5) (5.1,1.9) (5.9,2.1) (5.6,1.8) (5.8,2.2) (6.6,2.1)
    (5.5,1.8) (6.3,1.8) (5.8,1.6) (6.1,2.5) (5.1,2.0) (5.3,1.9)};
  \legend{Espèce A, Espèce B, Espèce C}
\end{axis}
\end{graphe}

On note un vecteur en gras, $\vect{x} \in \Rn$, et ses composantes
$x_1, x_2, \dots, x_n$. On les empile verticalement :
\[
  \vect{x} = \begin{pmatrix} x_1 \\ x_2 \\ \vdots \\ x_n \end{pmatrix}.
\]

\begin{intuition}
Voyez le vecteur $\vect{x}\in\R^2$ comme une \emph{flèche} partant de
l'origine et pointant vers le point $(x_1, x_2)$. Additionner deux vecteurs,
c'est mettre les flèches bout à bout.
\end{intuition}

\begin{schema}{Le vecteur $\vect{x} = (3,4)$ vu comme une flèche. Sa
longueur se lit directement avec le théorème de Pythagore : les pointillés
gris forment un triangle rectangle de côtés $3$ et $4$, dont l'hypoténuse
— la flèche — mesure $5$.}[vecteur-fleche]
  \begin{scope}[x=1cm, y=1cm]
    \axesxy{5.2}{5.2}
    \draw[construction] (3,0) -- (3,4);
    \draw[construction] (0,4) -- (3,4);
    \draw[vecteur] (0,0) -- (3,4);
    \node[point] at (3,4) {};
    \node[etiquette, below] at (3,-0.12) {$x_1 = 3$};
    \node[etiquette, left]  at (-0.12,4) {$x_2 = 4$};
    \node[font=\small, above left=-2pt] at (1.5,2) {$\vect{x}$};
    \node[annotation bulle] (mblong) at (4.3,1.6)
         {longueur\\ $\norm{\vect{x}}_2 = 5$};
    \draw[trait annotation] (mblong) to[bend right=18] (2.1,2.5);
  \end{scope}
\end{schema}

\begin{definition}[Norme euclidienne][norme-eucl]
La \emph{longueur} d'un vecteur $\vect{x}\in\Rn$, appelée \emph{norme
euclidienne}, est le nombre
\[
  \norm{\vect{x}}_2 = \sqrt{x_1^2 + x_2^2 + \dots + x_n^2}
                    = \sqrt{\smash{\textstyle\sum_{i=1}^{n} x_i^2}}.
\]
C'est exactement le théorème de Pythagore, généralisé à $n$ dimensions.
\end{definition}

\begin{exemple}[Calcul d'une norme]
Pour $\vect{x} = (3, 4)$, on a $\norm{\vect{x}}_2 = \sqrt{9 + 16}
= \sqrt{25} = 5$. La flèche mesure 5~unités de long.
\end{exemple}

\begin{piege}
La norme n'est \textbf{pas} la somme des composantes ! On élève au carré,
on additionne, \emph{puis} on prend la racine. Oublier la racine est
l'erreur numéro~un des débutants.
\end{piege}

\section{Une matrice transforme les vecteurs}

Une matrice $A \in \Mmn{m}{n}$ est un tableau de nombres à $m$ lignes et
$n$ colonnes. Son vrai rôle : \textbf{transformer} un vecteur en un autre,
par le produit $\vect{y} = A\vect{x}$.

\begin{aretenir}
Multiplier par une matrice, c'est appliquer une transformation
géométrique : une rotation, un étirement, une projection\dots{} Comprendre
\emph{quelle} transformation cache une matrice, c'est comprendre l'algèbre
linéaire.
\end{aretenir}

\begin{enia}
Une couche de réseau de neurones n'est rien d'autre que cela : une
multiplication par une matrice, $\vect{y} = W\vect{x} + \vect{b}$, suivie
d'une petite fonction appliquée nombre par nombre. Les coefficients de $W$
sont précisément ce que le réseau \emph{apprend}. Entraîner un modèle,
c'est chercher les bons nombres à mettre dans des matrices — et rien de
plus mystérieux que cela.
\end{enia}

\begin{schema}{Un petit réseau de neurones. Chaque trait porte un
coefficient : ce sont exactement les cases de la matrice $W$. Un réseau
qui « apprend », c'est un réseau qui ajuste ces nombres.}[reseau-couche]
  \reseauneurones[Entrée $\vect{x}$][Couche cachée][Sortie $\vect{y}$]{3}{4}{2}
\end{schema}

\begin{theoreme}[Existence et unicité de la solution][inversible]
Soit $A \in \Mn$ une matrice carrée. Le système linéaire $A\vect{x} =
\vect{b}$ admet une \emph{unique} solution pour tout $\vect{b}\in\Rn$ si et
seulement si $A$ est \emph{inversible}, c'est-à-dire $\det(A) \neq 0$.
Dans ce cas, $\vect{x} = \inv{A}\,\vect{b}$.
\end{theoreme}

\begin{proof}
Si $A$ est inversible, poser $\vect{x} = \inv{A}\vect{b}$ donne
$A\vect{x} = A\inv{A}\vect{b} = \identity\vect{b} = \vect{b}$, d'où
l'existence ; et si $A\vect{x} = A\vect{x}'$ alors, en multipliant à gauche
par $\inv{A}$, on obtient $\vect{x} = \vect{x}'$, d'où l'unicité.
Réciproquement, si $A$ n'est pas inversible, son noyau $\Ker A$ contient un
vecteur non nul et la solution, quand elle existe, n'est jamais unique.
\end{proof}

\begin{remarque}
En pratique, on ne calcule \emph{presque jamais} $\inv{A}$ pour résoudre
$A\vect{x}=\vect{b}$ : c'est coûteux et numériquement instable. On préfère
des méthodes directes (l'élimination de Gauss, au chapitre suivant) ou
itératives. Voir aussi~\textcite{golubvanloan2013}.
\end{remarque}

\section{Passons au code}

Voici comment manipuler tout cela en Python avec \co{python}{numpy}, la
bibliothèque de référence en calcul scientifique~\parencite{harris2020numpy}.

\begin{code}{python}[Norme et résolution d'un système linéaire]
import numpy as np

# Un vecteur et sa norme euclidienne
x = np.array([3.0, 4.0])
print(np.linalg.norm(x))          # -> 5.0

# Résoudre A x = b SANS inverser A (recommandé)
A = np.array([[2.0, 1.0],
              [1.0, 3.0]])
b = np.array([1.0, 2.0])

x = np.linalg.solve(A, b)         # méthode directe, stable
print("solution :", x)
print("residu   :", np.linalg.norm(A @ x - b))
\end{code}

Le même calcul en Julia, très prisé en recherche numérique :

\begin{code}{julia}[La même résolution en Julia]
using LinearAlgebra

A = [2.0 1.0; 1.0 3.0]
b = [1.0, 2.0]

x = A \ b            # l'opérateur « backslash » résout le système
println("solution : ", x)
println("norme    : ", norm(x))
\end{code}

\begin{plusloin}
Le nombre $\cond(A) = \norm{A}\,\norm{\inv{A}}$, appelé
\emph{conditionnement}, mesure à quel point la solution est sensible aux
erreurs d'arrondi. Un grand conditionnement annonce des ennuis numériques.
Nous y reviendrons en détail.
\end{plusloin}

\section*{Exercices}
\addcontentsline{toc}{section}{Exercices}

\begin{exercice}[Norme et distance]
Calculez $\norm{\vect{x}}_2$ pour $\vect{x} = (1, 2, 2)$, puis la distance
entre $\vect{x}$ et $\vect{y} = (1, 0, 0)$.
\end{exercice}

\begin{solution}
$\norm{\vect{x}}_2 = \sqrt{1 + 4 + 4} = 3$. La distance vaut
$\norm{\vect{x}-\vect{y}}_2 = \norm{(0,2,2)}_2 = \sqrt{8} = 2\sqrt{2}$.
\end{solution}

\begin{exercice}[Inversibilité]
La matrice $A = \begin{psmallmatrix} 1 & 2 \\ 2 & 4 \end{psmallmatrix}$
est-elle inversible ? Justifiez, puis dites ce que cela implique pour le
système $A\vect{x} = \vect{b}$.
\end{exercice}

\begin{solution}
$\det(A) = 1\cdot 4 - 2\cdot 2 = 0$ : la matrice \textbf{n'est pas}
inversible (la 2\ieme{} ligne est le double de la 1\iere{}). D'après le
\cref{thm:inversible}, le système n'a pas de solution unique : soit aucune
solution, soit une infinité.
\end{solution}
